\documentclass[12pt,a4paper]{article} \usepackage[utf8]{inputenc}
 \usepackage{amsmath, amssymb} 
 \usepackage{geometry} 
 \usepackage{hyperref} 
 \usepackage{listings}
  \usepackage{xcolor} 
  \geometry{margin=2cm}

\definecolor{codegray}{gray}{0.9} \lstset{ backgroundcolor=\color{codegray}, basicstyle=\ttfamily\small, breaklines=true, frame=single }

\begin{document}

\title{Exercícios Resolvidos e Programação MATLAB/Python\\large Formação em Sistemas de Comunicação Avançados} \author{} \date{} \maketitle

\tableofcontents \newpage

%%%%%%%%%%%%%%%%%%%%%%%%%%%%%%%%%%%%%%%%%%%%%%%%%%%%%%%%%%%%%%% \section{Exercícios Resolvidos – Fundamentos} %%%%%%%%%%%%%%%%%%%%%%%%%%%%%%%%%%%%%%%%%%%%%%%%%%%%%%%%%%%%%%%

\subsection{Exercício 1 — Cálculo de $E_b/N_0$} Um sistema QPSK transmite a 4 Mbps com potência de 10 mW. O ruído tem densidade espectral $N_0 = 10^{-9}$ W/Hz.

\textbf{Resolução:} \begin{align*} E_b &= \frac{P}{R_b} = \frac{10 \times 10^{-3}}{4 \times 10^6} = 2.5 \times 10^{-9} \ \frac{E_b}{N_0} &= 2.5 \ 10\log_{10}(2.5) &\approx 3.98,\text{dB} \end{align*}

\textbf{Resposta:} $E_b/N_0 \approx 4$ dB.

\subsection{Exercício 2 — BER teórico de QPSK} Para $E_b/N_0 = 10$ dB: \begin{align*} BER &= Q\left(\sqrt{2\frac{E_b}{N_0}}\right) = Q(4.472) \approx 3.87 \times 10^{-6} \end{align*}

\textbf{Resposta:} $BER \approx 3.9 \times 10^{-6}$.

%%%%%%%%%%%%%%%%%%%%%%%%%%%%%%%%%%%%%%%%%%%%%%%%%%%%%%%%%%%%%%% \section{Exercícios Resolvidos – Modulação} %%%%%%%%%%%%%%%%%%%%%%%%%%%%%%%%%%%%%%%%%%%%%%%%%%%%%%%%%%%%%%%

\subsection{Exercício 3 — Energia média de 16-QAM} \begin{align*} E_s &= \frac{2}{3}(M-1) = 10 \end{align*}

\textbf{Resposta:} $E_s = 10$.

\subsection{Exercício 4 — Efeito do ruído em 64-QAM} A constelação dispersa radialmente, reduzindo a separação entre símbolos e aumentando o BER.

%%%%%%%%%%%%%%%%%%%%%%%%%%%%%%%%%%%%%%%%%%%%%%%%%%%%%%%%%%%%%%% \section{Exercícios Resolvidos – SDR e Canais} %%%%%%%%%%%%%%%%%%%%%%%%%%%%%%%%%%%%%%%%%%%%%%%%%%%%%%%%%%%%%%%

\subsection{Exercício 5 — Passo de quantização} ADC de 12 bits, gama de $[-1,1]$ V: \begin{align*} \Delta &= \frac{2}{4096} = 4.88 \times 10^{-4},\text{V} \end{align*}

\textbf{Resposta:} $0.488$ mV.

\subsection{Exercício 6 — Doppler em UAV} \begin{align*} f_d &= \frac{v}{c} f_c = 1600,\text{Hz} \end{align*}

\textbf{Resposta:} $f_d \approx 1.6$ kHz.

%%%%%%%%%%%%%%%%%%%%%%%%%%%%%%%%%%%%%%%%%%%%%%%%%%%%%%%%%%%%%%% \section{Exercícios Resolvidos – Algoritmos PHY} %%%%%%%%%%%%%%%%%%%%%%%%%%%%%%%%%%%%%%%%%%%%%%%%%%%%%%%%%%%%%%%

\subsection{Exercício 7 — Equalizador LMS} \begin{align*} w(n+1) &= w(n) + \mu e(n)x(n) = 0.502 \end{align*}

\textbf{Resposta:} $w(n+1)=0.502$.

\subsection{Exercício 8 — Costas Loop} Adequado para QPSK porque resolve ambiguidades de fase sem piloto.

%%%%%%%%%%%%%%%%%%%%%%%%%%%%%%%%%%%%%%%%%%%%%%%%%%%%%%%%%%%%%%% \section{Exercícios Resolvidos – Simulação} %%%%%%%%%%%%%%%%%%%%%%%%%%%%%%%%%%%%%%%%%%%%%%%%%%%%%%%%%%%%%%%

\subsection{Exercício 9 — BER em simulação} \begin{align*} BER &= \frac{120}{100000} = 1.2 \times 10^{-3} \end{align*}

%%%%%%%%%%%%%%%%%%%%%%%%%%%%%%%%%%%%%%%%%%%%%%%%%%%%%%%%%%%%%%% \section{Exercícios Resolvidos – Fixed-Point} %%%%%%%%%%%%%%%%%%%%%%%%%%%%%%%%%%%%%%%%%%%%%%%%%%%%%%%%%%%%%%%

\subsection{Exercício 10 — Conversão para Q1.7} \begin{align*} 0.375 \times 128 &= 48 \end{align*}

\textbf{Resposta:} 00110000.

%%%%%%%%%%%%%%%%%%%%%%%%%%%%%%%%%%%%%%%%%%%%%%%%%%%%%%%%%%%%%%% \newpage \section{Exercícios de Programação MATLAB e Python} %%%%%%%%%%%%%%%%%%%%%%%%%%%%%%%%%%%%%%%%%%%%%%%%%%%%%%%%%%%%%%%

\subsection{Constelações QPSK e 16-QAM}

\subsubsection*{MATLAB} \begin{lstlisting}[language=Matlab] M = 4; data = randi([0 M-1], 1000, 1); symbols_qpsk = pskmod(data, M, pi/4); figure; plot(symbols_qpsk, 'o'); grid on;

M = 16; data = randi([0 M-1], 1000, 1); symbols_qam = qammod(data, M, 'UnitAveragePower', true); figure; plot(symbols_qam, 'o'); grid on; \end{lstlisting}

\subsubsection*{Python} \begin{lstlisting}[language=Python] import numpy as np import matplotlib.pyplot as plt

M = 4 data = np.random.randint(0, M, 1000) symbols_qpsk = np.exp(1j * (np.pi/4 + data * np.pi/2)) plt.scatter(symbols_qpsk.real, symbols_qpsk.imag) plt.show()

M = 16 data = np.random.randint(0, M, 1000) I = 2*(data % 4) - 3 Q = 2*(data // 4) - 3 symbols_qam = (I + 1j*Q) / np.sqrt(10) plt.scatter(symbols_qam.real, symbols_qam.imag) plt.show() \end{lstlisting}

%%%%%%%%%%%%%%%%%%%%%%%%%%%%%%%%%%%%%%%%%%%%%%%%%%%%%%%%%%%%%%% \subsection{Simulação AWGN e BER}

\subsubsection*{MATLAB} \begin{lstlisting}[language=Matlab] N = 1e5; M = 4; EbN0_dB = 8; data = randi([0 M-1], N, 1); tx = pskmod(data, M, pi/4); rx = awgn(tx, EbN0_dB, 'measured'); data_hat = pskdemod(rx, M, pi/4); BER = sum(data ~= data_hat) / N \end{lstlisting}

\subsubsection*{Python} \begin{lstlisting}[language=Python] N = int(1e5); M = 4; EbN0_dB = 8

data = np.random.randint(0, M, N) tx = np.exp(1j * (np.pi/4 + data * np.pi/2)) EbN0 = 10**(EbN0_dB/10) noise = (np.random.randn(N)+1jnp.random.randn(N))/np.sqrt(2EbN0) rx = tx + noise angles = np.angle(rx) data_hat = np.mod(np.round((angles - np.pi/4)/(np.pi/2)),4).astype(int) BER = np.mean(data != data_hat) \end{lstlisting}

%%%%%%%%%%%%%%%%%%%%%%%%%%%%%%%%%%%%%%%%%%%%%%%%%%%%%%%%%%%%%%% \subsection{Canal Rayleigh com Doppler}

\subsubsection*{MATLAB} \begin{lstlisting}[language=Matlab] fd = 200; chan = rayleighchan(1/1e6, fd); tx = pskmod(randi([0 3], 5000, 1), 4, pi/4); rx = filter(chan, tx); plot(chan.PathGains) \end{lstlisting}

\subsubsection*{Python} \begin{lstlisting}[language=Python] N = 5000; fd = 200; fs = 1e6 t = np.arange(N)/fs rayleigh = (np.random.randn(N)+1jnp.random.randn(N))/np.sqrt(2) doppler = np.exp(1j2np.pifd*t) h = rayleigh * doppler plt.plot(np.abs(h)) plt.show() \end{lstlisting}

%%%%%%%%%%%%%%%%%%%%%%%%%%%%%%%%%%%%%%%%%%%%%%%%%%%%%%%%%%%%%%% \subsection{AGC Simples}

\subsubsection*{MATLAB} \begin{lstlisting}[language=Matlab] rx = tx + 0.5randn(size(tx)); target = 1; gain = target ./ abs(rx); rx_agc = rx . gain; \end{lstlisting}

\subsubsection*{Python} \begin{lstlisting}[language=Python] rx = tx + 0.5*(np.random.randn(N)+1j*np.random.randn(N)) target = 1 gain = target / np.mean(np.abs(rx)) rx_agc = rx * gain \end{lstlisting}

%%%%%%%%%%%%%%%%%%%%%%%%%%%%%%%%%%%%%%%%%%%%%%%%%%%%%%%%%%%%%%% \subsection{Equalizador LMS}

\subsubsection*{MATLAB} \begin{lstlisting}[language=Matlab] mu = 0.01; w = 0; x = randn(1,1000); d = filter([1 0.5],1,x); for n=2:length(x) y(n)=wx(n); e(n)=d(n)-y(n); w=w+mue(n)*x(n); end \end{lstlisting}

\subsubsection*{Python} \begin{lstlisting}[language=Python] mu=0.01; w=0 x=np.random.randn(1000) d=np.convolve(x,[1,0.5],mode='same') y=np.zeros_like(x); e=np.zeros_like(x) for n in range(1,len(x)): y[n]=wx[n] e[n]=d[n]-y[n] w=w+mue[n]*x[n] \end{lstlisting}

%%%%%%%%%%%%%%%%%%%%%%%%%%%%%%%%%%%%%%%%%%%%%%%%%%%%%%%%%%%%%%% \subsection{Conversão para Fixed-Point}

\subsubsection*{MATLAB} \begin{lstlisting}[language=Matlab] function y = q17(x) y = round(x*2^7); y = max(min(y,127),-128); end \end{lstlisting}

\subsubsection*{Python} \begin{lstlisting}[language=Python] def q17(x): y=int(round(x*128)) return max(min(y,127),-128) \end{lstlisting}

%%%%%%%%%%%%%%%%%%%%%%%%%%%%%%%%%%%%%%%%%%%%%%%%%%%%%%%%%%%%%%% \end{document}